% Page 034 — page imprimée 30
% Contenu seul : le préambule et la pagination sont gérés par meyrueis.tex

« ...Pour les tourdres, je le regrette bien mais jusqu’ici il est impossible d’en trouver, la
surveillance est si sévère que les paysans n’osent pas s’aventurer. On parle de 6 mois de
prison et de 5 ou 600 F. d’amende. Si, un peu plus tard, je peux en trouver quelques uns soyez
assuré que je ferai mon possible mais je ne puis rien vous promettre. » M. Gasc à Requista

\medskip

\textbf{La guerre de 1914-18}

Sur le plan politique les menaces s’accumulent et l’on entend des bruits de botte. Fin 1912
Albert est convoqué pour une période d’instruction militaire. 18 mois plus tard ce sera la
guerre, et Albert sera tué lors de la bataille de Soissons.

Son père comptait sur lui pour reprendre et relancer la chapellerie après la disparition de Paul,
l’aîné.

\medskip

Que lui était-il donc arrivé ?

Dans les années 90 Paul était tombé amoureux fou d’une de ses cousines germaines Emma,
fille de Louis . Elle l’aimait bien et l’appréciait comme cousin mais pas plus et elle ne lui
laissa aucun espoir.

Paul en fut terriblement affecté et lorsque l’occasion s’en présenta (?) il partit pour le Chili
pour oublier et probablement aussi pour prendre des contacts pour un éventuel commerce de
chapeaux avec ce pays. Muni d’un petit pécule il embarqua sur un vapeur pour Santiago et...
« on n’eut plus jamais aucune nouvelle de lui. » Est-il mort de maladie ? A-t-il été assassiné ?
Nul ne le sait. C’est là la tradition Loupiac. A Cambonéral chez les Verdier la mémoire est
restée plus vive : « Il nous a toujours été raconté que Paul s’était embauché dans une
chapellerie au Chili. Il avait été happé par une machine et en serait mort. » On peut penser
que Léonie sœur cadette de Paul Couderc donnera à son fils en souvenir de son frère disparu
le même prénom afin de ne pas l’oublier...

Les deux fils disparus , Auguste se retrouvait seul avec ses deux filles ( Marie infirme et
célibataire et Marthe qui mourra d’un cancer à Montpellier en 1925, célibataire elle aussi. )
pour continuer la chapellerie mais il avait 74 ans. La guerre continuait avec son cortège de
morts. Découragé et fatigué il arrêta la fabrication et ferma la chapellerie à peu près au même
moment que les établissements Veyrier pourtant beaucoup plus importants.

On ne parlera plus désormais des chapeliers de Meyrueis.

La guerre cette « maudite guerre » comme avait dit son frère Louis avait disloqué beaucoup
de familles de la campagne, laissant des veuves et des filles qui ne se marièrent pas car des
millions de jeunes hommes avaient été massacrés. Dès lors avec la modernisation, les petites
filatures, les tissages et les chapelleries de notre ville de Meyrueis étaient condamnés. Et qui
peut comprendre la souffrance des survivants ?

.
